В курсовом проекте реализован согласованный вычислительный контур для работы с нечёткими величинами на основе U(a,b). C++-ядро задаёт математическую модель, реализует хранение α-срезов, интервальную арифметику, сравнение, генерацию и кластеризацию. C API делает ядро доступным из Go, а Go-сервер обеспечивает проверку запросов, синхронное состояние и HTTP-интеграцию. Веб-интерфейс и CEF-оболочка используют один и тот же серверный контракт.

Поставленные задачи выполнены. Проверки \texttt{fuzzy\_tests} подтверждают корректность базовых инвариантов: вложенность срезов, границы операций, воспроизводимость генератора, корректность отбора и устойчивое разделение синтетических кластеров. Контрольные измерения показывают миллисекундное время обработки наборов до 5000 интервалов в Release-сборке, при этом пороговая кластеризация закономерно требует больше времени, чем k-means++ с фиксированным k.

Практическая область применения модуля включает учебные эксперименты, анализ интервалов с неопределёнными границами и прототипирование алгоритмов нечёткого кластерного анализа. При переходе к промышленной нагрузке потребуются расширенные бенчмарки, анализ устойчивости к выбору α и δ, а также отдельное хранилище состояния.

---